%!TEX root = ../main.tex
\section{Related Work}
\label{sec:related}

\subsection{The ComParE 2017 Cold Sub-Challenge}

The ComParE 2017 organizers evaluated 6,373-dimensional acoustic
functionals, bag-of-audio-words representations and CNN/LSTM systems, both
alone and in late fusion~\cite{schuller2017interspeech}. The official 0.710
hidden-Test UAR is the majority-vote fusion of all baseline families, not the
standalone openSMILE/SVM result. Contemporary submissions likewise emphasized
complementarity. Huckvale and Beke compared voice quality, vowel, modulation
and spectral-distribution features~\cite{huckvale2017cold}; Cai \etal used
constant-Q and Gammatone spectra with end-to-end networks~\cite{cai2017cold}.
These findings motivated our compact energy/pause and constant-Q branches.

Huckvale and Beke also give the clearest contemporary statement of the
speaker-confounding risk on this corpus. Noting that the training set consists
of 10,000 extracts drawn from only 210 speakers, of whom a minority had a cold,
they warned that ``the classifier may just become a speaker recognizer for
those speakers in the training set that have a cold''~\cite{huckvale2017cold}.
Their results show the cost. Their best development configuration reached
70.97\% UAR, far above the 64.0\% development baseline, but the same system
returned 62.1\% on the hidden Test set, below the openSMILE/SVM baseline it had
beaten on Development. Our project reproduced this pattern before diagnosing
it, which is why we treat evaluation validity as a precondition for the
feature and model comparisons rather than as a closing check.

\subsection{Foundation models for speech and health audio}

Since 2017 the dominant representation has shifted from engineered functionals
to self-supervised transformers. wav2vec~2.0~\cite{baevski2020wav2vec2},
HuBERT~\cite{hsu2021hubert} and WavLM~\cite{chen2022wavlm} learn from
unlabelled audio and transfer to corpora far too small for end-to-end
training, which is precisely the URTIC regime. Their layers are not
interchangeable: layer-wise analyses show that lower layers retain acoustic
and phonatory detail while upper layers become progressively more
linguistic~\cite{pasad2021layer}. This matters for a task whose evidence is
phonatory rather than lexical, and it motivates the layer selection and
layer-weighting experiments we report in Section~\ref{sec:methods}.

A second line of work pretrains on health audio directly rather than on
speech in general. Google's HeAR is trained on more than 300 million short
clips of coughs, breathing and other human sounds, and reports
state-of-the-art results across 33 health acoustic
tasks~\cite{baur2024hear}. TRILLsson distils universal paralinguistic
representations from a large teacher model and targets exactly the
non-lexical properties a cold should perturb~\cite{shor2022trillsson}. Neither
had been evaluated on URTIC, and pretraining relevance is not the same as
demonstrated complementarity on a particular small corpus, so we test both
against the compact handcrafted branches rather than assuming they dominate.

\subsection{Speaker representations and the confound}

ECAPA-TDNN is the standard modern speaker-verification
embedding~\cite{desplanques2020ecapa}, and we use it to reconstruct the
speaker grouping the student release omits. This creates a tension the paper
has to confront. The same voice-quality dimensions that a cold perturbs,
phonation, breathiness and spectral tilt, are part of what a speaker embedding
encodes, so a representation built to identify talkers is neither a clean
confound detector nor an obviously invalid cold feature. We therefore validate
ECAPA externally on labelled corpora before using it as a grouping proxy
(Section~\ref{sec:data_eval}), and we keep it out of the cold classifiers
themselves.

\subsection{De-confounding and domain adaptation}

Where a nuisance factor is known, several families of methods aim to remove
it. Gradient-reversal domain-adversarial training suppresses a nuisance
variable by training the representation to defeat a
discriminator~\cite{ganin2015dann}, and margin-disparity discrepancy
tightened the underlying theory~\cite{zhang2019mdd}. Supervised contrastive
learning offers a different route, shaping the embedding by pulling together
same-label examples~\cite{khosla2020supcon}, which can be masked so that
positives come from different speakers. CORAL instead aligns source and
unlabelled-target feature statistics without touching
labels~\cite{sun2016coral}, a conservative option when the target is a hidden
Test partition. We evaluate the gradient-reversal and contrastive routes as
de-confounding candidates and diagonal CORAL as a domain-adaptation component
of one submitted system. A recurring finding, reported in
Section~\ref{sec:protocol}, is that each of these needs a matched control
before its effect can be attributed to the mechanism it claims.

\subsection{Sequential and spectral representations}

Two representation choices in our submitted systems predate the deep-learning
era but remain competitive at this sample size. The constant-Q
transform~\cite{brown1991cqt} uses geometrically spaced, constant-ratio
frequency bins, giving finer low-frequency resolution than a fixed-window
spectrogram and matching the region where phonation changes concentrate. Path
signatures summarize the ordered geometry of a multivariate sequence rather
than only its marginal statistics, and are established tools for sequential
learning~\cite{toth2020signatures}. Both retain information that
whole-utterance pooling discards, while adding far fewer parameters than a
foundation-model head.
